% Session 5: Subagents, Hooks, and Final Project
% Beamer Presentation --- Claude Code for Academic Researchers
\documentclass[aspectratio=169, 11pt]{beamer}

% ── Theme and Colors ──────────────────────────────────────────────
\usetheme{Madrid}
\usecolortheme{default}

\definecolor{anthropic}{RGB}{50, 50, 80}
\definecolor{accent}{RGB}{180, 90, 40}
\definecolor{lightgray}{RGB}{245, 245, 245}

\setbeamercolor{structure}{fg=anthropic}
\setbeamercolor{frametitle}{bg=anthropic, fg=white}
\setbeamercolor{title}{bg=anthropic, fg=white}
\setbeamercolor{block title}{bg=anthropic, fg=white}
\setbeamercolor{block body}{bg=lightgray}
\setbeamercolor{alerted text}{fg=accent}

\setbeamerfont{title}{size=\Large, series=\bfseries}
\setbeamerfont{frametitle}{size=\large, series=\bfseries}

% ── Packages ──────────────────────────────────────────────────────
\usepackage{graphicx}
\usepackage{booktabs}
\usepackage{listings}
\usepackage{fontenc}
\usepackage{hyperref}
\usepackage{tikz}
\usetikzlibrary{positioning}

\lstset{
  basicstyle=\ttfamily\small,
  backgroundcolor=\color{lightgray},
  frame=single,
  framerule=0pt,
  breaklines=true,
  columns=fullflexible,
  literate={_}{\_}1,
}

% ── Metadata ──────────────────────────────────────────────────────
\title{Session 5: Subagents, Hooks, and Final Project}
\subtitle{Spawning focused agents, event-driven automation, and an end-to-end research artifact}
\author{Brady Allardice}
\date{}
\institute{UPF \& IBEI}

% ══════════════════════════════════════════════════════════════════
\begin{document}

% ── Title ─────────────────────────────────────────────────────────
\begin{frame}[plain]
  \titlepage
\end{frame}

% ══════════════════════════════════════════════════════════════════
%  SUBAGENTS
% ══════════════════════════════════════════════════════════════════

\begin{frame}{Subagents}
  \centering
  \vspace{2em}
  {\Large A focused agent spawned for a bounded job ---\\[0.3em]
  the right tool for the parallelizable, scoped work\\[0.3em]
  that already shows up in your pipeline.}

  \vspace{2em}
  Concrete patterns, then named teams.
\end{frame}

\begin{frame}{What a Subagent Is}
  \textbf{Claude Code can spawn another Claude Code instance inside your current session --- a subagent.}

  \vspace{0.6em}
  \begin{itemize}
    \item It has its own context window
    \item It gets a scoped task and a scoped tool set
    \item It returns a summary to your main conversation --- not its raw trace
  \end{itemize}

  \vspace{0.8em}
  \textbf{What that buys you:}
  \begin{enumerate}
    \item \textbf{Parallelism} --- run three searches at once instead of one after the other
    \item \textbf{Context hygiene} --- a big messy exploration stays out of your main context budget
  \end{enumerate}

  \vspace{0.5em}
  \begin{block}{Think of it as}
    Delegating a bounded sub-task to a focused colleague who returns a clean, structured answer --- so your main conversation stays on the high-level work.
  \end{block}
\end{frame}

\begin{frame}{Where Subagents Shine}
  \textbf{Reach for a subagent when the subtask is:}
  \begin{itemize}
    \item \textbf{Parallelizable} --- independent work that doesn't need to be sequenced
    \item \textbf{Bounded} --- clear inputs, clear outputs, no back-and-forth
    \item \textbf{Independently verifiable} --- you can check the result without reading the full trace
  \end{itemize}

  \vspace{0.7em}
  \textbf{Concrete research use cases:}
  \begin{itemize}
    \item \textbf{Multi-database literature searches} --- one agent per database, all running in parallel
    \item \textbf{Coding many documents} against a fixed rubric
    \item \textbf{Parallel robustness checks} across independent specifications
    \item \textbf{Running the same sanity-check} script over many datasets
    \item \textbf{Exploring a codebase} --- send a subagent to map out the files, keep the summary in your main context
  \end{itemize}
\end{frame}

\begin{frame}{Using Subagents Well}
  \textbf{Three habits that make subagents pay off:}
  \begin{enumerate}
    \item \textbf{Brief them like a contractor.} Concrete deliverable, all the context they need up front --- they don't see your conversation history.
    \item \textbf{Ask for structured output} (bullet list, table, JSON). You'll be consuming the result; make it easy to read and verify.
    \item \textbf{Verify independently.} You get a summary, not the full trace. Spot-check the way you would a student's work.
  \end{enumerate}

  \vspace{0.8em}
  \begin{block}{The payoff}
    A literature sweep, robustness battery, or repo survey that would have taken an hour finishes in minutes --- and doesn't eat into your main conversation's context. That's what makes them worth reaching for.
  \end{block}
\end{frame}

% ══════════════════════════════════════════════════════════════════
%  NAMED SUBAGENTS / MULTI-AGENT TEAMS
% ══════════════════════════════════════════════════════════════════

\begin{frame}[fragile]{Building a Team of Specialized Agents}
  \textbf{Named subagents you define once and invoke on demand.} Each has its own scope, system prompt, and tool allowlist --- like a team of focused collaborators.

  \vspace{0.4em}
  \textbf{Define one:} a markdown file at \texttt{\textasciitilde/.claude/agents/<name>.md} (user-level) or \texttt{.claude/agents/<name>.md} (project-level).

  \vspace{0.2em}
  \begin{lstlisting}[basicstyle=\ttfamily\scriptsize]
---
name: lit-search
description: Use for multi-database literature searches on a topic.
tools: [mcp__zotero__*, WebSearch, WebFetch]
model: sonnet
---
You are a literature search specialist. Return a structured
bibliography with the key claim from each paper.
  \end{lstlisting}

  \vspace{0.2em}
  \textbf{Activate it:}
  \begin{itemize}\small
    \item Claude reads every agent's \texttt{description} at startup and spawns one when a task matches
    \item Or invoke it explicitly: \emph{``use the \texttt{lit-search} agent to\ldots''}
    \item Spawn several in a single message --- they run in parallel
  \end{itemize}

  \vspace{0.2em}
  \alert{\small Permissions work as usual: any tools the agent uses need to be allowed in \texttt{.claude/settings.json}.}
\end{frame}

\end{document}
